\subsection{Independent mathematical checks of editorial emendations}
\label{checks:1985}

This subsection checks corrections to established mathematics; it is not a second
research development selected from class C or D. The original notebook formulas
remain in the occurrence index. The labels \ours\ and \oursComputed\ mean that
the proof or computation was supplied for this companion, not that the underlying
mathematics is new. Exact code and output are in
\texttt{code/check\_mathematics.py} and its JSON output. A coefficient-algebra
implementation and direct application to functions are separate checks.

\begin{proposition}[The corrected Novikov--Veselov commutator]\ours\
Let \(D=\partial_z\), \(E=\partial_{\bar z}\), and let the smooth coefficients
\(V,u,v\) satisfy \(Eu=3DV\), \(Dv=3EV\). Set
\[
 L=DE+V,\qquad P=D^3+E^3+uD+vE,\qquad B=Du+Ev.
\]
Here a coefficient without a derivative to its right denotes multiplication.
Then
\[
 [P,L]+BL=M_{D^3V+E^3V+D(uV)+E(vV)}.
\]
For real \(V\), one may put \(v=\bar u\). Consequently the right hand side
defines the stated fixed-energy evolution, with the boundary or normalization
conditions on \(u\) still to be specified.
\end{proposition}
\begin{proof}
Leibniz's rule gives the following coefficients, before imposing the auxiliary
equations:
\[
\begin{array}{c|l}
\text{derivative} & \text{coefficient in }[P,L]+BL\\\hline
D^2&3DV-Eu\\
E^2&3EV-Dv\\
DE&B-Du-Ev\\
D&3D^2V-DEu\\
E&3E^2V-DEv\\
1&D^3V+E^3V+uDV+vEV+BV.
\end{array}
\]
All positive-order coefficients vanish; the scalar coefficient is the one
claimed. This calculation also fixes the meaning of the notebook expression
\((D u)V\): it must be \(D(uV)\), not \((Du)V\) alone. Independently, the code
applies the two operators successively to two different test polynomials, for
each of two polynomial potentials, constructing \(u,v\) by integration of the
auxiliary equations. Every residual is exactly zero. This second check tests the
implementation; the displayed algebra proves the identity for all smooth
coefficients. The hierarchy itself is the prior result of
\citet[\S II, equations (4)--(6)]{veselovnovikov1984explicit}.
\end{proof}

\begin{proposition}[An explicit counterexample to the unrestricted no-Lax wording]\ours\
There is a smooth real potential \(U(x,t)\), independent of \(y\), for which
\[
 L=-\partial_x^2-\partial_y^2+U,\qquad
 A=-4\partial_x^3+6U\partial_x+3U_x
\]
satisfy \(L_t=[A,L]\), but the potentials do not evolve by Euclidean motions.
The example is not a genuinely two-dimensional decaying scattering potential;
that extra hypothesis is precisely absent from the universal notebook claim.
\end{proposition}
\begin{proof}
Put
\[
 r=e^{2x-8t},\quad s=e^{4x-64t},\quad
 T=1+r+s+\tfrac19rs,\quad U=-2\partial_x^2\log T.
\]
Since \(T>0\) for real \(x,t\), the formula is smooth and real for all such
parameters. To verify the equation without relying on a remembered soliton
formula, use the four exponential terms of \(T\), with pairs
\((k,\omega)=(0,0),(2,-8),(4,-64),(6,-72)\) and coefficients
\(1,1,1,1/9\). In the bilinear expression
\((D_x^4+D_xD_t)T\mathbin{\cdot}T\), each ordered pair contributes
\((k_i-k_j)^4+(k_i-k_j)(\omega_i-\omega_j)\).
Every unordered pair gives zero except \((0,3)\) and \((1,2)\). They have the
same exponential \(rs\) and weighted coefficients \(864/9=96\) and \(-96\),
respectively, so they cancel. Writing \(q=\log T\), the bilinear identity is
\(q_{xt}+q_{xxxx}+6q_{xx}^2=0\); differentiation gives
\(U_t=6UU_x-U_{xxx}\).

A direct Leibniz expansion gives
\([A,L]=M_{6UU_x-U_{xxx}}\), since \(A\) commutes with
\(\partial_y^2\). Thus the Lax equation holds. To check that this is not a
time-dependent translation, compare \(U_t/U_x\) at \(t=0\), first with
\((r,s)=(1,1)\), then with \((r,s)=(2,4)\). The exact ratios are $-92/3$ and $-89524/4861$, respectively, and both denominators $U_x$ are nonzero. These values were also obtained independently from moment and cumulant formulas for the four exponential terms. A translation would give a
ratio depending on \(t\) alone. A continuous Euclidean motion of a nonconstant
potential independent of \(y\), remaining independent of \(y\), cannot rotate
its direction of variation, so it reduces to this excluded translation case.
The code independently verifies the rational PDE identity in \(r,s\) and the
formal commutator by two implementations. This counterexample does not assert
that the intended genuinely two-dimensional obstruction has been identified.
\end{proof}

\begin{lemma}[Leading Baker--Akhiezer coefficients]\ours\
Use the notebook's exponential conventions and suppose locally that
\[
 \psi=e^{ik_1z}(1+\xi_1/k_1+O(k_1^{-2})),\qquad
 \psi=c_2e^{ik_2\bar z}(1+O(k_2^{-1})),\quad c_2\ne0.
\]
For an operator \(L=DE+aE+b\) annihilating \(\psi\), cancellation of the
leading terms requires
\(a=-D\log c_2\), \(b=-iE\xi_1\).
\end{lemma}
\begin{proof}
After dividing by \(e^{ik_1z}\), the constant term of \(DE\psi\) is
\(iE\xi_1\); the term \(aE\psi\) begins at order \(k_1^{-1}\).
At the other puncture the coefficient of \(ik_2e^{ik_2\bar z}\) is
\(Dc_2+ac_2\). These two coefficients give the formulas. The recorded
operator \(DE+A D+B\) instead has an uncancelled \(ik_1 A\) term at the first
puncture. Symbolic differentiation of two truncated Laurent expansions gives
the same cancellations independently. Existence and uniqueness from spectral
data are supplied by the cited theorem of
\citet[\S3, Lemmas 1--2]{dubrovinkrichevernovikov1976schrodinger}; this lemma
only determines the necessary coefficients in the convention used here.
\end{proof}

\begin{proposition}[The hyperelliptic defect in products of differentials]\ours\
Let \(X\) be a smooth compact connected complex hyperelliptic curve of genus \(g\ge2\).
The multiplication map
\(\operatorname{Sym}^2H^0(X,K_X)\to H^0(X,K_X^2)\) has image dimension
\(2g-1\) and cokernel dimension \(g-2\).
\end{proposition}
\begin{proof}
Choose a coordinate on the quotient so that infinity is unbranched, and write
\(y^2=f(x)\) with \(f\) squarefree of degree \(2g+2\). The one-forms
\(x^i dx/y\), \(0\le i\le g-1\), are a basis, by local regularity and
Riemann--Roch. Their products span the independent forms
\[
 x^k\frac{dx^2}{y^2},\quad 0\le k\le 2g-2.
\]
The complementary forms are
\[
 x^j\frac{dx^2}{y},\quad 0\le j\le g-3,
\]
with the second list empty when \(g=2\). At a branch point use
\(x-a=t^2\), \(y=t\) times a nonvanishing analytic factor: the two basic
quadratic differentials have orders zero and one, respectively. At either
point at infinity put \(s=1/x\). Their orders are \(2g-2-k\) and
\(g-3-j\), so all listed forms are holomorphic. The two families have opposite
parity under the hyperelliptic involution; inside each family their rational
monomials are independent. There are \(3g-3=h^0(K_X^2)\) forms in total,
so they form a basis. This proves both dimensions for every \(g\ge2\).
Independently the code checks the orders, product exponents and dimensions for
\(g=2,\ldots,9\), explicitly including the exceptional zero-cokernel case
\(g=2\). The general nonhyperelliptic surjectivity statement is Max Noether's
theorem, recalled in \citet[\S3.1]{matonevolpato2013determinantal}.
\end{proof}

\begin{lemma}[The sign of the torus Beltrami wedge]\ours\
Write \(y=\Im\tau>0\), and use the deformation
\(\tau_\eta=(\tau+\eta\bar\tau/y)/(1+\eta/y)\). At \(\eta=0\),
\[
 d\tau=-2i\,d\eta,\qquad
 d\bar\eta\wedge d\eta=-\tfrac14d\tau\wedge d\bar\tau.
\]
\end{lemma}
\begin{proof}
Differentiate the quotient with the base \(\tau\) fixed; the derivative is
\((\bar\tau-\tau)/y=-2i\). Conjugation and antisymmetry give the wedge
identity. For an independent real-coordinate check, \(\eta=p+iq\) gives
\((\delta\Re\tau,\delta\Im\tau)=(2q,-2p)\), with Jacobian determinant four.
The coordinate identities \(d\bar\eta\wedge d\eta=2i\,dp\wedge dq\) and
\(d\tau\wedge d\bar\tau=-2i\,d\Re\tau\wedge d\Im\tau\) give exactly
the same minus sign. These are tangent identities at the base point.
\end{proof}

\begin{proposition}[The finite-dimensional determinant metric fixes the reciprocal]\ours\
For a linear map \(T:E_1\to E_2\) of finite-dimensional Hermitian spaces, put
\(\lambda_T=\det\ker T\otimes(\det\operatorname{coker}T)^{-1}\).
Under the determinant-of-complex isomorphism
\(\lambda_T\simeq\det E_1\otimes(\det E_2)^{-1}\), the norm induced from the
fixed spaces is the cohomology norm divided by the product of the nonzero
singular values of \(T\).
\end{proposition}
\begin{proof}
Take orthonormal bases of \(\ker T\), its orthogonal complement, the image,
and its orthogonal complement, with \(Te_i=\sigma_i f_i\), \(\sigma_i>0\).
Complete the kernel basis by the \(e_i\), and the cokernel basis by the
\(Te_i\), in a fixed determinant sign convention. The corresponding element
of \(\det E_1\otimes(\det E_2)^{-1}\) contains the scalar
\(\prod_i\sigma_i^{-1}\). The remaining bases are orthonormal, which proves
the formula; signs have modulus one. If \(T\) is invertible, the canonical
acyclic generator therefore maps to \((\det T)^{-1}\), whereas \(\det T\)
belongs to the dual line. Empty products give the rank-zero case.

The positive eigenvalues of \(T^*T\) are \(\sigma_i^2\), so the squared norm
factor is \((\det{}'T^*T)^{-1}\). Two unrelated rational diagonal examples
and the empty example are checked in the code. Replacing the product by a
zeta determinant yields the direction of the spectral factor on the
notebook's line; the analytic existence and smoothness theorem remains
\citet[\S3]{quillen1985determinants}, who uses its dual.
\end{proof}

\begin{example}[Anomaly coefficients and modular-weight check]\oursComputed\
For \(x=c_1(K_{X/S})\), the degree-two term of
\[
 e^{jx}\left(1-\frac{x}{2}+\frac{x^2}{12}\right)
\]
is \((6j^2-6j+1)x^2/12\). This is also obtained by multiplying the three
terms by hand: \(j^2/2-j/2+1/12\). Thus \(c_1=1\), \(c_2=13\), and
\(c_2-13c_1=0\). This is only the coefficient calculation; the analytic
family-index theorem is cited in the occurrence index.

For a modular substitution, the genus-one factors
\(d^2\tau\), \((\Im\tau)^{-14}\), and \(|\eta(\tau)|^{-48}\) have
absolute weights \(-4,28,-24\), whose sum is zero. The corresponding
\(q\)-density has \((-\log|q|)^{-14}\). Separately, the torus formula
\(C y^2|\eta(\tau)|^4\) for \(|dz|^2\), read in
\citet[\S5.0.1, equation (5.1)]{kokotov2007polyhedral}, becomes
\(C y|\eta(\tau)|^4\) in the unit-area metric: its Laplacian is multiplied
by \(y\), and \(\zeta_\Delta(0)=-1\). The latter follows from the flat
two-dimensional heat expansion after subtracting the one-dimensional
constant eigenspace. Area conventions cannot be changed in only one factor
of the measure.
\end{example}

\paragraph{Scope of the second checks.}
The Novikov--Veselov and KdV identities were rederived in coefficient form and
by direct application. The KdV example also has a separate bilinear proof.
The determinant reciprocal was checked by both a one-dimensional acyclic map
and the singular-value argument. The Beltrami sign was checked in complex and
real coordinates. The hyperelliptic dimension result was checked by local
orders and Riemann--Roch, with finite examples used only to catch boundary
mistakes. The stacking development has an independent exact facet-enumeration
check and a separate proof audit recorded with its source files. No finite
computation is presented as a proof of an unbounded universality assertion.
